% Part III -- Reusable and Reliable Programs
% Chapter 8: Errors and Exceptions

\h{Errors and Exceptions}
\label{ch:08}

Errors are normal while learning to program. The useful skill is not avoiding
every error. It is reading what Python reports, locating the cause, and deciding
whether the program should correct the problem or explain it to the user.

\begin{NoteBox}
\textbf{Prerequisites.} You should understand type conversion, conditionals,
loops, and functions with parameters and return values.
\end{NoteBox}

\begin{tcolorbox}[title=Learning outcomes,colframe=orange,colback=white,arc=2mm]
After completing this chapter, you should be able to:
\begin{itemize}
    \item distinguish syntax, runtime, and logic errors;
    \item read the final line of a traceback;
    \item catch a specific exception with \c{try} and \c{except};
    \item use \c{else} and \c{finally} when they clarify the flow;
    \item validate repeated input without hiding unrelated errors; and
    \item raise a clear exception inside a function.
\end{itemize}
\end{tcolorbox}

\hh{Three kinds of problems}

Program problems often fall into three broad groups:

\begin{itemize}
    \item \textbf{Syntax error}: Python cannot understand the written structure.
    \item \textbf{Runtime exception}: Python begins running but an operation fails.
    \item \textbf{Logic error}: the program runs but produces the wrong result.
\end{itemize}

This intentionally incorrect instruction has an unclosed string:

\begin{lstlisting}[
    language=text,
    style=block,
    caption={Text showing a syntax mistake; do not run it as correct code},
    label={c08_syntax_example},
    captionpos=b
]
print("Lobby)
\end{lstlisting}

A syntax problem must be corrected before the program can run. A runtime
exception can sometimes be anticipated and handled. A logic error requires
checking the calculation or decision rule.

\begin{ExplainBox}{An exception is an event and a transfer of control}
When an operation raises an exception, Python stops the normal sequence at
that operation and searches outward for a matching handler. Statements later
in the abandoned block do not run. A matching \c{except} block is therefore an
alternate path through the program, not a way to erase the original event.
If no matching handler exists, the program stops and prints a traceback.
\end{ExplainBox}

Decide how to respond from the cause, not from a desire to hide red text:
\begin{itemize}
    \item fix a syntax or programming defect rather than catching it;
    \item handle an expected external failure when the program has a meaningful
          recovery, such as asking again after invalid input;
    \item raise an exception when a function cannot honor its contract;
    \item allow an unexpected exception to remain visible during development
          so its traceback preserves evidence about the defect.
\end{itemize}

A handler is justified only if it can state what failed, why that failure was
anticipated, and what safe next state the program will enter.

\hh{Read a traceback from the bottom}

When an exception is not handled, Python prints a traceback. Begin with its
last line. It names the exception and describes the immediate problem. Then
look upward for the line in your file.

\begin{lstlisting}[
    language=python,
    style=block,
    caption={An unhandled ValueError example},
    label={c08_runtime_error},
    captionpos=b
]
width_text = "three"
width_m = float(width_text)
print(width_m)
\end{lstlisting}

The final traceback line includes \c{ValueError}. The failed operation is the
attempt to convert \c{"three"} to a floating-point number.

\bookfigure[0.9\textwidth]{Figures/Generated/ch08_traceback_anatomy.pdf}
{Read a traceback from the exception message upward to the relevant line in your file.}
{fig:ch08-traceback-anatomy}

\begin{TryBox}{Create an error on purpose}
Run the conversion example exactly as written. Identify the exception type,
the message, and the line in your file. Then replace \c{"three"} with
\c{"3.0"} and explain why the same operation now succeeds.
\end{TryBox}

\hh{Catch a specific exception}

Place only the operation that may fail inside \c{try}. Catch the exception you
expect by name.

\begin{lstlisting}[
    language=python,
    style=block,
    caption={Handle an expected conversion error},
    label={c08_basic_try_except},
    captionpos=b
]
width_text = "three"

try:
    width_m = float(width_text)
except ValueError:
    print("Enter a number such as 3.5.")
\end{lstlisting}

The program does not crash. It gives a message that helps the user correct the
input.

\begin{SyntaxBox}{Exception handling separates possible paths}
The \c{try} block contains the operation that may fail. A matching
\c{except ValueError:} block runs only for that exception. An optional
\c{else} block is the success path, and an optional \c{finally} block runs
after either outcome. Catch the narrowest exception that you can explain.
\end{SyntaxBox}

\hh{Keep the try block narrow}

A narrow \c{try} block makes it clear which operation is protected.

\begin{lstlisting}[
    language=python,
    style=block,
    caption={Separate conversion from the success path},
    label={c08_narrow_try},
    captionpos=b
]
width_text = "6.0"
length_m = 8.0

try:
    width_m = float(width_text)
except ValueError:
    print("Width must be numeric.")
else:
    area_m2 = width_m * length_m
    print(f"Area: {area_m2:.2f} m2")
\end{lstlisting}

If the calculation in \c{else} contains a different bug, Python reports it
instead of incorrectly treating it as a conversion problem.

\hh{Handle different exceptions separately}

Different failures often require different messages.

\begin{lstlisting}[
    language=python,
    style=block,
    caption={Give each expected failure a useful message},
    label={c08_multiple_exceptions},
    captionpos=b
]
total_text = "24.0"
count_text = "0"

try:
    total = float(total_text)
    count = int(count_text)
    average = total / count
except ValueError:
    print("Total and count must be numbers.")
except ZeroDivisionError:
    print("Count must be greater than zero.")
else:
    print(f"Average: {average:.2f}")
\end{lstlisting}

Catch specific exceptions before broader ones. Avoid a bare \c{except:}
because it can hide mistakes that should be fixed.

\hh{Use the exception message carefully}

Add \c{as error} when the original message is useful for diagnosis.

\begin{lstlisting}[
    language=python,
    style=block,
    caption={Capture the original exception message},
    label={c08_exception_object},
    captionpos=b
]
value_text = "4O"  # note: that's the letter O, not the digit zero

try:
    value = float(value_text)
except ValueError as error:
    print(f"Conversion failed: {error}")
\end{lstlisting}

System messages are helpful for a programmer but may be unclear to a beginner.
A user-facing program can provide a plain explanation and record technical
details separately when needed.

\hh{Use else for successful completion}

The \c{else} block runs only when the \c{try} block completes without an
exception.

\begin{lstlisting}[
    language=python,
    style=block,
    caption={Run a success action in else},
    label={c08_else_block},
    captionpos=b
]
height_text = "3.2"

try:
    height_m = float(height_text)
except ValueError:
    print("Height must be numeric.")
else:
    print(f"Accepted height: {height_m:.2f} m")
\end{lstlisting}

Use \c{else} when it makes the success path easier to distinguish. It is not
required after every \c{try} block.

\hh{Use finally for an action that must run}

The \c{finally} block runs whether the operation succeeds or fails.

\begin{lstlisting}[
    language=python,
    style=block,
    caption={Run a final status action},
    label={c08_finally_block},
    captionpos=b
]
value_text = "not a number"

try:
    value = float(value_text)
    print(value)
except ValueError:
    print("Conversion failed.")
finally:
    print("Conversion attempt finished.")
\end{lstlisting}

Chapter 10 uses with blocks to close files automatically. That is usually
clearer than manual cleanup in \c{finally}.

\bookfigure[0.84\textwidth]{Figures/Generated/ch08_exception_flow.pdf}
{The try block leads to a matching except block after an expected failure or to else after success; finally runs on either path.}
{fig:ch08-exception-flow}

The four clauses are not interchangeable. Keep risky operations narrow in
\c{try}, recover only from errors you understand in \c{except}, keep success
work in \c{else}, and reserve \c{finally} for actions that truly must run.

\hh{Repeat until input is valid}

A \c{while} loop can ask again after an expected conversion failure.

\begin{lstlisting}[
    language=python,
    style=block,
    caption={Ask again until a positive number is entered},
    label={c08_input_validation_loop},
    captionpos=b
]
while True:
    value_text = input("Enter a positive width in meters: ")

    try:
        width_m = float(value_text)
    except ValueError:
        print("Use a number such as 3.5.")
        continue

    if width_m <= 0:
        print("Width must be greater than zero.")
        continue

    break

print(f"Accepted: {width_m:.2f} m")
\end{lstlisting}

Conversion failure and value validation are separate. The exception handles
text that is not numeric. The \c{if} statement handles a numeric value that is
outside the accepted range.

\hh{Raise an exception in a function}

A function can use \c{raise} to report that an argument violates its rules.

\begin{lstlisting}[
    language=python,
    style=block,
    caption={Raise and handle a clear ValueError},
    label={c08_raise_value_error},
    captionpos=b
]
def rectangle_area_m2(width_m, length_m):
    if width_m <= 0 or length_m <= 0:
        raise ValueError("width and length must be positive")

    return width_m * length_m

try:
    print(rectangle_area_m2(-2.0, 5.0))
except ValueError as error:
    print(f"Cannot calculate area: {error}")
\end{lstlisting}

The function reports the invalid argument. The caller decides whether to show
a message, ask again, or stop.

\hh{Do not use exceptions for ordinary decisions}

Use a condition when a situation is expected and easy to test.

\begin{lstlisting}[
    language=python,
    style=block,
    caption={Use a condition for an expected empty list},
    label={c08_condition_before_operation},
    captionpos=b
]
values = []

if len(values) == 0:
    print("No values are available.")
else:
    print(sum(values) / len(values))
\end{lstlisting}

Use exception handling when an operation may fail for reasons that are awkward
or unreliable to predict in advance, such as converting user text or opening a
file.

\hh{Common exception types}

\begin{center}
\begin{tabular}{ll}
\textbf{Exception} & \textbf{Typical cause} \\
\c{ValueError} & A value has the right type but unsuitable content \\
\c{TypeError} & An operation receives an unsuitable type \\
\c{ZeroDivisionError} & A calculation divides by zero \\
\c{NameError} & A name has not been defined \\
\c{IndexError} & A list position does not exist \\
\c{KeyError} & A dictionary key does not exist \\
\c{FileNotFoundError} & A requested file cannot be found \\
\end{tabular}
\end{center}

The exception name helps locate the category of problem. The surrounding code
still needs to be examined.

\hh{Common mistakes}

\begin{itemize}
    \item \textbf{Catching every error}: broad handling hides programming mistakes.
    \item \textbf{Putting too much in try}: the source of failure becomes unclear.
    \item \textbf{Ignoring the exception}: an empty handler leaves the user without guidance.
    \item \textbf{Using the wrong exception type}: the handler never runs.
    \item \textbf{Continuing with an undefined result}: success code belongs in \c{else}
          or after a guaranteed assignment.
    \item \textbf{Raising a vague message}: state which value or rule is invalid.
\end{itemize}

\hh{Practice ladder}

\hhh{Level 0 -- Read and predict}

\begin{lstlisting}[
    language=python,
    style=block,
    caption={Predict which block runs},
    label={c08_practice_level0},
    captionpos=b
]
try:
    value = int("12")
except ValueError:
    print("Invalid")
else:
    print(value + 1)
\end{lstlisting}

\hhh{Level 1 -- Modify}

Change the text to \c{"twelve"}. Run the program and explain why the other
block executes.

\hhh{Level 2 -- Complete}

Complete the handler by replacing \c{...} with a statement that displays a
useful message.

\begin{lstlisting}[
    language=python,
    style=block,
    caption={Complete a division handler},
    label={c08_practice_level2},
    captionpos=b
]
numerator = 20
denominator = 0

try:
    print(numerator / denominator)
except ZeroDivisionError:
    ...
\end{lstlisting}

\hhh{Level 3 -- Apply}

Write \c{parse_floor_count(text)}. Convert the text to an integer. Raise a
\c{ValueError} if the result is less than 1. Call it inside \c{try}/\c{except}.

\textbf{Hints}
\begin{itemize}
    \item \textbf{Hint 1}: \c{int(text)} already raises \c{ValueError} for unsuitable text.
    \item \textbf{Hint 2}: Use an \c{if} statement for a converted value below 1.
    \item \textbf{Hint 3}: Return the accepted count.
\end{itemize}

\hhh{Level 4 -- Challenge}

Use a loop to ask for a floor count until \c{parse_floor_count()} returns a
valid value. Display the reason after each rejected entry.

\hh{Practice solutions}

\textbf{Level 0:} The output is \c{13}.

\textbf{Level 2}

\begin{lstlisting}[
    language=python,
    style=block,
    caption={One solution to Level 2},
    label={c08_solution_level2},
    captionpos=b
]
numerator = 20
denominator = 0

try:
    print(numerator / denominator)
except ZeroDivisionError:
    print("Denominator must not be zero.")
\end{lstlisting}

\textbf{Level 3}

\begin{lstlisting}[
    language=python,
    style=block,
    caption={One solution to Level 3},
    label={c08_solution_level3},
    captionpos=b
]
def parse_floor_count(text):
    count = int(text)
    if count < 1:
        raise ValueError("floor count must be at least 1")
    return count

try:
    floor_count = parse_floor_count("0")
except ValueError as error:
    print(f"Invalid floor count: {error}")
else:
    print(f"Accepted: {floor_count}")
\end{lstlisting}

\textbf{Level 4}

\begin{lstlisting}[
    language=python,
    style=block,
    caption={One solution to Level 4},
    label={c08_solution_level4},
    captionpos=b
]
def parse_floor_count(text):
    count = int(text)
    if count < 1:
        raise ValueError("floor count must be at least 1")
    return count

while True:
    text = input("Floor count: ")
    try:
        floor_count = parse_floor_count(text)
    except ValueError as error:
        print(f"Try again: {error}")
        continue
    break

print(f"Accepted: {floor_count}")
\end{lstlisting}

\hh{Chapter checkpoint}

\begin{enumerate}
    \item What are the three broad kinds of program problems described here?
    \item Where should you begin when reading a traceback?
    \item Why should an \c{except} clause name a specific exception?
    \item When does \c{else} run after a \c{try} block?
    \item When does \c{finally} run?
    \item Why should a \c{try} block be narrow?
    \item What does \c{raise} do?
    \item Should a numeric value below an accepted limit require conversion handling?
    \item Why is a bare \c{except:} risky?
\end{enumerate}

\hh{Checkpoint answers}

\begin{enumerate}
    \item Syntax errors, runtime exceptions, and logic errors.
    \item At the final line, then trace upward to your code.
    \item It handles only the expected failure and leaves other bugs visible.
    \item Only when the \c{try} block finishes without an exception.
    \item Whether the \c{try} block succeeds or fails.
    \item It makes the failing operation easier to identify.
    \item It creates an exception and ends the current function or operation.
    \item No. Convert first, then use a condition to check the range.
    \item It can catch and hide errors that were not anticipated.
\end{enumerate}

\begin{tcolorbox}[title=Chapter summary,colframe=orange,colback=white,arc=2mm]
You can classify errors, read a traceback, catch specific exceptions, separate
success and cleanup paths, validate repeated input, and raise clear errors from
functions. The next chapter explains why two variable names may share the same
mutable object.
\end{tcolorbox}

